\chapter{Non-24-Hour Sleep-Wake Rhythm Disorder}
\index{Non-24-Hour Sleep-Wake Rhythm Disorder}
\label{sec:Non-24-Hour Sleep-Wake Rhythm Disorder}

\section{Overview}
Non-24-hour sleep-wake rhythm disorder (N24) is a circadian rhythm sleep-wake disorder characterized by a sleep-wake cycle that is desynchronized from the 24-hour day, typically running at a period longer than 24 hours (free-running), with each day sleep onset and wake time drifting progressively later by 30--60 minutes. It occurs in 50--80\% of blind individuals (due to lack of photic entrainment of the SCN) and rarely in sighted individuals. This condition is among the most common mental health disorders worldwide. It affects people across all age groups, socioeconomic backgrounds, and geographic regions.
\section{Causes}
This condition happens because of absence of light perception to entrain the circadian clock.     
\section{Risk Factors}

\begin{itemize}[noitemsep]
  \item Genetic predisposition and family history
  \item Advancing age
  \item Lifestyle factors (diet, exercise, smoking, alcohol)
  \item Environmental exposures
  \item Underlying medical conditions
  \item Medication use
\end{itemize}
\section{Signs and Symptoms}

\subsection{Common symptoms}
\begin{itemize}[noitemsep]
  \item You may experience a progressive daily delay in sleep-wake timing cycling through the 24-hour day. 
  \end{itemize}

\subsection{Emergency warning signs}
\begin{itemize}[noitemsep]
  \item Severe or worsening symptoms of non-24-hour sleep-wake rhythm disorder require immediate medical evaluation
\end{itemize}
\section{Progression and Stages}
The condition may progress through different stages of severity over time. Early stages often involve mild or subtle symptoms that may be overlooked. Moderate stages involve more noticeable symptoms affecting daily function. Advanced stages may involve significant impairment and complications.
\section{Complications}
\begin{itemize}[noitemsep]
  \item Disease progression leading to worsening symptoms
  \item Secondary infections or injuries
  \item Side effects of treatment
  \item Reduced quality of life
  \item Emotional and psychological impact
\end{itemize}
\section{Diagnosis}
\begin{itemize}[noitemsep]
  \item To make the diagnosis, doctors look for at least 14 days of actigraphy or sleep diaries demonstrating a circadian period other than 24 hours. Management in blind individuals includes melatonin (0.5--10 mg at desired bedtime) and tasimelteon (approved for N24 in blind people)
  \item In sighted individuals, Treatment options include morning light therapy and evening melatonin.
  \item Comprehensive medical history and physical examination
  \item Laboratory tests to assess organ function
  \item Imaging studies to visualize affected structures
  \item Specialized testing depending on the affected system
  \item Consultation with relevant specialists
\end{itemize}
\section{Prevention}
\begin{itemize}[noitemsep]
  \item Maintain a healthy lifestyle with balanced diet and exercise
  \item Avoid known risk factors
  \item Attend regular medical check-ups for early detection
  \item Follow recommended screening guidelines
\end{itemize}
\section{Treatment Options}
Medications to manage symptoms and address underlying mechanisms Lifestyle modifications and self-care measures Physical therapy and rehabilitation Surgical intervention when indicated Regular monitoring and follow-up care Patient education and support services

\begin{itemize}[noitemsep]
  \item Medications to manage symptoms and address underlying mechanisms
  \item Lifestyle modifications and self-care measures
  \item Physical therapy and rehabilitation
  \item Surgical intervention when indicated
  \item Regular monitoring and follow-up care
\end{itemize}
\section{Living with It}
\begin{itemize}[noitemsep]
  \item Follow your treatment plan as prescribed
  \item Maintain regular follow-up with your healthcare provider
  \item Make necessary lifestyle adjustments
  \item Seek support from family, friends, or support groups
  \item Stay informed about your condition
\end{itemize}
\section{Outlook}
outlook is chronic but manageable. The outlook varies depending on the specific condition, severity at diagnosis, and response to treatment. Early intervention generally leads to better outcomes. Many conditions can be effectively managed with appropriate treatment. Regular follow-up care is important for monitoring and treatment adjustment.
